\scanpage{12}
% L1
Comme ceci est vrai pour tt $y$, on trouve

% L2
\[
S\geq\mathrm{Sup}\bigl((U'\cup U'')\times V_y\bigr)
 =(U'\cup U'')\times\mathrm{Sup}\,V_y=(U'\cup U'')\times V
\]

% L3
contrairement au caractère maximal de $U'$, OK.

% L4
\uline{Cas de produits \deleted{\unsure{infinis|infin.}} d'une famille quelconque de topos}

% L5–8
On trouve que si tous les facteurs sauf \inserted{au plus} un seul sont tels que tt pt ait un système fond. de voisinages (pas néc. ouverts) quasi-compacts, et si tous sauf au plus un seul, sont \deleted{\unsure{munis}} eux-mêmes quasi-compacts, alors

% L9
\[
\mathrm{Top}\bigl(\prod_{i\in I}X_i\bigr)\longrightarrow\prod_{i\,(\mathrm{top})}^{\unsure{(2)|(b)}}\mathrm{Top}(X_i)
\]

% L10
est une équivalence de topos.

% L11–14
\deleted{C'est essentiellement le $=$ argument, travaillant sur la cat. \unsure{du} produit restreint $\prod' O_{X_i}$ (familles \deleted{\unsure{de|des}} $(U_i)_{i\in I}$ $U_i\in O_{X_i}$ pour tt $i$, $U_i=X_i$ (\inserted{plus} \deleted{\unsure{objet}} grand objet de $O_{X_i}$) pour presque tout $i$) et la ``topologie produit'' sur ce produit.}

% L15–17
Dans le cas où $I$ est fini, c'est un corollaire immédiat de ce qui précède. Dans le cas général, on a $I=\varinjlim_\alpha I_\alpha$ ($I_\alpha$ parties finies de $I$) donc

% L18–19
\[
\begin{aligned}
\prod_{i\,\mathrm{top}}\mathrm{Top}(X_i)
 &\cong\mathop{\varprojlim_{\mathrm{top}}}_{\alpha}\bigl(\prod_{i\in I_\alpha}\mathrm{Top}(X_i)\bigr)
 \cong\mathop{\varprojlim_{\mathrm{top}}}_{\alpha}\bigl(\mathrm{Top}(\underbrace{\prod_{i\in I_\alpha}X_i}_{Z_\alpha}\unsure{)}\bigr)\\
 &\cong\bigl(\underbrace{\varinjlim_\alpha O_{Z_\alpha}}_{O'_Z}\bigr)^{\sim}
\end{aligned}
\]

% L20–26
$O'_Z$ est la catégorie des ouverts de $Z$ \inserted{$=\prod X_i$} qui sont images inverses d'ouverts de produits finis $\prod_{i\in I_\alpha}X_i=Z_\alpha$. \deleted{\unsure{Ils}} On \deleted{\unsure{peut|une}} a dessus la topologie \inserted{$\pi$ des sites} $\varinjlim$, qui \deleted{\unsure{définit}} est utilisée dans la formule précédente, mais aussi la top. plus fine $\pi'$ provenant de la top de $Z$. Si tous les \deleted{\unsure{Z}} \unsure{$X_i$|$i$} sauf un nb fini sont quasi-compacts [les autres étant ce que ils veulent] on trouve par un argument tt analogue \inserted{à dessus} que les top. sont

% foot
% archive sheet number: 11
